\documentclass[11pt]{article}
\usepackage[a4paper,margin=25mm]{geometry}
\usepackage{amsmath,amssymb,amsthm,hyperref}
\newtheorem{theorem}{Theorem}
\newtheorem{lemma}{Lemma}
\newcommand{\PP}{\mathcal P}
\title{Local two-variable completion for a polynomial density}
\author{Complete proof; three independent mathematical reviews completed}
\date{30 September 2026}
\begin{document}\maketitle
The purpose of this note is to replace one positive polynomial density
and one uniform density, on a single interval, by two rational densities
on alternating pure cells. The cell masses have prescribed denominators.
No bound on the denominators of higher moments or cell endpoints is asserted.
All constants below can depend on fixed positive bounds $a,b$ for the input
density, and on no other input.

\begin{theorem}\label{main}
Let $w\in\mathbb Q[t]$ be a probability density on $[0,1]$, of degree $d$,
with $0<a\le w\le b$. Put $R=m+d+2$ and
$\mu=\int_0^1t w(t)\,dt$. There is a constant $C$ such that, for every
pair of integers
\[
 K\ge CR^4,\qquad L\ge C(KR)^{160},\qquad KL\mu\in\mathbb Z,
\]
there are rational piecewise polynomial probability densities $f_A,f_B$
with the following properties.
\begin{enumerate}
\item The support of $f_B$ consists of $K$ disjoint rational boxes,
each of mass $1/K$. The support of $f_A$ is their complement, and its
mass on every complementary gap is a positive integer divided by $L$.
\item On every nonempty pure cell, the conditional density rescaled to
$[0,1]$ is bounded above and below by positive constants depending only
on $a,b$. Each polynomial piece has degree at most $m$, except that
the $B$ pieces have degree at most one.
\item If $A,B$ have these densities, then all their ordered polynomial
moments of total degree at most $m$ agree with those of independent
variables $A_0\sim\operatorname{Unif}[0,1]$ and $B_0$ of density $w$.
In particular adding $m$ independent uniforms gives the same complete
order probabilities as the original pair and those uniforms.
\end{enumerate}
\end{theorem}

\begin{lemma}[Separated quadrature for a bounded polynomial density]\label{quad}
For $D\ge d+1$ and every integer $K\ge C(D+d+1)^4$, there is a real
equal-weight degree-$D$ quadrature for $w$ whose nodes satisfy
\[
 |b_q-F_w^{-1}((q-1/2)/K)|\le c/(8K),
\]
where $c>0$ depends only on $a,b$. Its endpoint and consecutive gaps
are bounded below by $c/K$.
\end{lemma}
\begin{proof}
Write $Q=F_w^{-1}$. The assumptions and Markov's inequality give
$|Q'|\le C$ and $|Q''|\le C(d+1)^2$.
Use the same fixed-right-inverse Newton argument as the separated-rule
lemma in
\path{../rational_designs/asymptotic/two_discrete_step_compatibility.tex},
with the midpoint quantile nodes $t_q=Q((q-1/2)/K)$.
For completeness, let $\phi_j$ be the uniform orthonormal shifted
Legendre polynomials and $f_j(t)=\int_0^t\phi_{j-1}(z)\,dz$, $1\le j\le D$.
The midpoint error for $f_j\circ Q$, using the displayed bounds on $Q$,
gives a residual vector of norm at most $C(D+d+1)^3/K^2$.
The derivative Gram matrix is
\[
 KAA^T=K^{-1}\sum_q\Phi(t_q)\Phi(t_q)^T,
 \qquad A_{jq}=\phi_{j-1}(t_q)/K.
\]
Its limiting matrix is $\int\Phi\Phi^T w\ge a I$; the rectangle
estimate gives operator error at most $C(D+d+1)^4/K$.
Thus its fixed right inverse has $2$-to-$\infty$ norm at most $CD$.
Within sup-norm radius $\rho$, the derivative variation has norm at most
$CD^3\rho$ from sup-norm inputs to Euclidean outputs. The Newton map
therefore contracts on its affine correction space in radius
$C(D+d+1)^4/K^2$. Increasing the constant in the lower bound on $K$
makes this radius smaller than $c/(8K)$ and the contraction constant
smaller than $1/2$. The quantile gaps are themselves at least $1/(bK)$,
which proves separation and interiority.
\end{proof}

\begin{proof}[Proof of Theorem~\ref{main}]
Choose an odd integer $h>d+m$, with $h\le d+m+2$, and set
$H(t)=P_h(2t-1)$. Let $D=2h+2m+2=O(R)$ and choose an exact rule
$b^0$ from Lemma~\ref{quad}. Increasing $C$ permits every $K\ge CR^4$.
Set
\[
 \eta=c_*(KR)^{-20},\qquad r_q^0=b_q^0-\eta H(b_q^0),
\]
where $c_*>0$ is a sufficiently small fixed constant.
The map $t\mapsto t-\eta H(t)$ has derivative at least $3/4$ and,
since $h$ is odd, maps the endpoints to $\eta$ and $1-\eta$.
Hence the $r_q^0$ are interior and separated by at least $c/K$.
Legendre orthogonality applies to the polynomial $t^j w(t)$:
its degree is less than $h$ for $0\le j\le m$. Exactness of the rule
therefore gives
\begin{align}
 K^{-1}\sum_q(b_q^0)^j&=\int t^j w(t)\,dt &&(0\le j\le m+1),\label{ord}\\
 K^{-1}\sum_q r_q^0(b_q^0)^j&=\int t^{j+1}w(t)\,dt &&(0\le j\le m).\label{rank}
\end{align}
In particular $\sum r_q^0=K\mu$.
Because $KL\mu$ is an integer, balanced rounding supplies integers
$a_q$ with $|a_q/L-r_q^0|\le1/L$ and $\sum a_q=KL\mu$.
Put $r_q=a_q/L$. They remain strictly increasing and interior.

We now correct $b$ with $r$ fixed. Use the $m+1$ ordinary moment
constraints in antiderivatives of $\phi_0,\ldots,\phi_m$, and the $m$
weighted constraints in $r_q f_j(b_q)-H_j(b_q)$, where
$H_j(t)=\int_0^t z\phi_{j-1}(z)\,dz$.
Their targets are the corresponding integrals against $w$.
The baseline derivative columns, multiplied by $K$, evaluate
\[
 \phi_0,\ldots,\phi_m,\quad
 -\eta H\phi_0,\ldots,-\eta H\phi_{m-1}.
\]
Their products have degree at most $D$, so their empirical Gram matrix
is exactly their continuous Gram matrix with weight $w$.
The Legendre multiplication argument in the cited step-compatibility
note applies also when $h=O(R)$ and $m<h$: Gaussian quadrature at the
zeros of $H$ bounds $\|p\|_2$ by $CR\|p+\eta Hq\|_2$, and its
sublevel estimate gives $\|Hq\|_2\ge cR^{-4}\|q\|_2$.
Consequently the weighted Gram matrix is at least
$c\eta^2R^{-10}I$, and the fixed right inverse has norm at most
$C\eta^{-2}R^{11}$.
The rounding residual has Euclidean norm at most $C\sqrt m/L$.
Exactly as in the cited proof, the fixed-inverse map contracts in
radius
\[
 \rho=C\eta^{-2}R^{12}/L.
\]
Indeed its contraction bound is at most
$C\eta^{-2}R^{11}(R^3\rho+R/L)$, which is less than $1/2$ under
the stated $L$ bound. The same bound makes $\rho\le\eta/2$, $1/L\le\eta/2$, and
$\rho\le c/(16K)$. We obtain separated nodes $b$ satisfying
\eqref{ord}--\eqref{rank} with $r$ in place of $r^0$, and
$\max|r_q-b_q|\le2\eta$.
The resulting fixed-$r$ Jacobian still has full row rank, because the
Newton map was a contraction on the chosen correction space.

The positive-completion proof in
\path{../independent_directions/two_exceptional_positive_completion.tex}
now applies without change. Its target for $A$ remains Lebesgue
measure; only the identities for the moments of $B$ have changed.
Here are the exact compatibility and positivity points.
For $S_j(a)=K^{-1}\sum_{b_q>a}b_q^j$ and
\[
 V=\operatorname{span}\{a^j:j\le m\}
   +\operatorname{span}\{a^rS_s(a):r+s\le m\},
\]
the cell-constant kernel is $S=\operatorname{span}(1,S_0,\ldots,S_m)$.
Giving the gaps masses $r_{q+1}-r_q$, including $r_0=0,r_{K+1}=1$,
annihilates the residual on this kernel, since
\[
 \sum_q r_q b_q^j/K=\sum_q b_q^{j+1}/K=\int_0^1 S_j(a)\,da.
\]
The quantitative quotient estimate in that proof gives a positive
completion as soon as
$\max|r_q-b_q|\le c\Delta^{5/2}/(m+1)^3$, where $\Delta$ is the
minimum gap. Since $\Delta\ge c/K$, our choice of $\eta$ meets this
condition with arbitrary fixed slack. The resulting point-node
completion density on the gaps lies between $1/2$ and $3/2$.

Finally apply the rational smoothing and rational Riesz-completion
arguments in \path{two_polynomial_exceptional_dice.tex}.
Replace each $b_q$ by a short rational box of mass $1/K$, with a
rational affine conditional density. The fixed-$r$ Jacobian just
proved has full rank, so its slope matrix is $\varepsilon/3$ times
that Jacobian plus $O(\varepsilon^3)$; solving a fixed rational minor
preserves exactly \eqref{ord}--\eqref{rank} with positive slopes.
All targets are rational because $w$ is rational polynomial.
The complementary-gap Riesz problem has the same step kernel, remains
uniformly close to the positive point-node completion, and has rational
Gram matrix and right side. It therefore gives a rational polynomial
density $f_A$ between $1/4$ and $7/4$. Its gap masses are still the
prescribed positive integers divided by $L$. The conditional density
bounds follow upon normalizing the cells.
For $r+s\le m$ its ordered moments are
\[
 \mathbb E[A^rB^s\mathbf1_{A<B}]
 =\frac1{r+1}\int b^{r+s+1}w(b)\,db,
\]
exactly those of the target pair. The ordinary moments and the opposite
chamber follow from independence and the same constraints. This proves
the theorem.
\end{proof}
\end{document}
